% This document was generated by an LLM.

\documentclass{essay}

\title{What Justifies Splitting $dy/dx$?}
\subtitle{On the Meaning of Differentials in Exact Equations}
\author{}
\date{}

\begin{document}

\maketitle

\section{The Puzzle}

Consider the chain-rule identity for $F(x,y)$ along a curve $y = y(x)$:
\[
\frac{dF}{dx} = \frac{\partial F}{\partial x}\bigg|_{(x,y(x))} + \frac{\partial F}{\partial y}\bigg|_{(x,y(x))}\frac{dy}{dx}.
\]
This equation underlies the theory of exact differential equations $M\,dx + N\,dy = 0$. Yet its use raises an old discomfort. Elementary calculus insists, correctly, that $dy/dx$ is \emph{not} a fraction but a single symbol denoting a limit. Later, when solving equations of the form $N\,dy/dx = -M$, we are casually told to ``multiply through by $dx$'' and write $M\,dx + N\,dy = 0$, treating $dy/dx$ exactly like a fraction after all. What justifies this? And underneath the notation, what actually \emph{is} a differential?

\section{The Chain Rule Is Not Fraction-Splitting}

It is worth first noticing that the identity above does not, by itself, split anything. The multivariable chain rule
\[
\frac{dF}{dx} = \frac{\partial F}{\partial x}\bigg|_{(x,y(x))} + \frac{\partial F}{\partial y}\bigg|_{(x,y(x))}\frac{dy}{dx}
\]
treats $dy/dx$ as one indivisible object: the ordinary derivative of $y$ with respect to $x$, valid because $y$ is assumed to be a differentiable function of $x$ along the curve in question. No fraction is decomposed here; this is simply differentiation of a composite function.

The real manipulation occurs one step later. Starting from
\[
\frac{dF}{dx} = M + N\frac{dy}{dx} = 0,
\]
we rearrange to $dy/dx = -M/N$, and then ``cross-multiply'' to obtain
\[
M\,dx + N\,dy = 0.
\]
This last move is the one that needs justifying.

\section{Two Objects, One Notation}

The resolution is that calculus silently uses the symbol $dy/dx$ for two different mathematical objects, and the discomfort comes from not distinguishing them.

\subsection*{The derivative}
This is
\[
\frac{dy}{dx} := \lim_{\Delta x \to 0} \frac{\Delta y}{\Delta x}, \qquad \Delta y := f(x + \Delta x) - f(x),
\]
a single limiting value. As a limit, it is genuinely not a fraction: both $\Delta y$ and $\Delta x$ vanish in the limiting process, and writing the result as ``$dy/dx$'' is a historical notation, not a literal quotient of two numbers called $dy$ and $dx$. This is precisely what elementary courses warn against, and the warning is correct.

\subsection*{The differential}
Once the derivative $f'(x)$ exists, one may \emph{define} two new quantities. Let $dx$ be an arbitrary nonzero real number---an independent variable, not required to be small---and define
\[
dy := f'(x)\,dx.
\]
Now $dy$ and $dx$ are genuine real numbers related by an ordinary algebraic equation. Dividing both sides by $dx$ gives $dy/dx = f'(x)$, trivially, by cancellation, precisely because $dy$ was constructed that way. This $dy/dx$ \emph{is} a fraction, and cross-multiplying it is ordinary algebra.

The two objects are built to agree numerically, which is why they share notation. But logically, the differential is manufactured \emph{after} the derivative exists, specifically so that their ratio reproduces it. The warning ``it's not a fraction'' applies to the derivative. The cross-multiplication trick used in solving $M\,dx + N\,dy = 0$ is legitimate because at that point the discussion has quietly shifted to the differential, without the shift being announced.

What is still missing is a precise statement of \emph{what kind of object} each side of this shift actually is---not just the equation it satisfies, but its type: its domain, its codomain, and the properties that come with it. Section~5 makes this explicit.

\section{What a Differential Actually Is}

There are three standard rigorous accounts of the differential, of increasing sophistication.

\subsection*{(a) Linear approximation}
Here $dy = f'(x)\,dx$ is defined as the \emph{linear part} of the true increment $\Delta y = f(x + dx) - f(x)$. Precisely,
\[
\Delta y = dy + \varepsilon(dx)\cdot dx, \qquad \varepsilon(dx) \to 0 \text{ as } dx \to 0.
\]
So $dy$ is not the actual change in $y$, but the tangent line's prediction of that change. This is the sense in which $dx$ and $dy$ are genuine numbers once a point and an increment are fixed.

\subsection*{(b) Linear map}
The modern, coordinate-free definition treats the differential of $F(x,y)$ at a point as the linear map (see Appendix~\ref{app:background} if the terms \emph{linear map} and \emph{tangent space} below are unfamiliar)
\[
dF_{(x,y)}(h,k) = \frac{\partial F}{\partial x}(x,y)\,h + \frac{\partial F}{\partial y}(x,y)\,k,
\]
where $dx$ and $dy$ are themselves defined as the coordinate projections $dx(h,k) = h$ and $dy(h,k) = k$---literally the differentials of the coordinate functions, which is why the names are consistent rather than coincidental. Then
\[
dF = \frac{\partial F}{\partial x}\,dx + \frac{\partial F}{\partial y}\,dy
\]
is an exact identity between linear functionals: no infinitesimals, no hand-waving, only linear algebra. It expresses $dF$ in the dual basis $\{dx, dy\}$. Evaluating both sides on the tangent vector $(h, y'(x)h)$ of a curve $y = y(x)$ and dividing by $h$ recovers the chain-rule identity directly---not by splitting a fraction, but by evaluating a linear map on a vector.

\subsection*{(c) Infinitesimal}
Non-standard analysis takes $dx$ to be an actual nonzero infinitesimal hyperreal number, and defines $dy := f(x + dx) - f(x)$ as the true increment, computed within the hyperreals. Then $dy/dx$ is a literal quotient of hyperreal numbers, infinitely close to $f'(x)$:
\[
\operatorname{st}\!\left(\frac{dy}{dx}\right) = f'(x).
\]
This is arguably the most faithful formalization of how Leibniz and Euler actually reasoned: $dx$ really is a number, just not a standard real one.

\section{The Type Signatures of Differential Objects}

The three accounts above all write down a symbol $dx$ and a symbol $dy$, and all three relate them by some version of $dy = f'(x)\,dx$. But $dx$ in framework (a) and $dx$ in framework (b) are not the same kind of mathematical object merely dressed differently---they have different domains and codomains outright. Stating each object's type, in the sense of a function signature (inputs in, output out), is what finally pins down what a differential \emph{is} in a given framework, as opposed to what equation it satisfies.

\subsection*{Objects common to all three frameworks}

Before the frameworks diverge, some notation means the same thing throughout.

\begin{itemize}[leftmargin=*]
  \item $x$: a real number, a point of the domain where $f$ is being differentiated. (In \S6 below, $x$ is reused for the first coordinate function $x : \R^2 \to \R$; context disambiguates the two uses.)
  \item $f$: a function $f : I \to \R$, where $I \subseteq \R$ is an open interval, differentiable at the point $x$ in question.
  \item $f'(x)$: for fixed $x$, a real number---the value of the derivative at that point. Letting $x$ vary, $f' : I \to \R$ is itself a function, related to $f$ by the usual linearity, product, and chain rules.
\end{itemize}

Everything that follows is about what \emph{additional} object, named $dx$ or $dy$, each framework builds on top of this common ground.

\subsection*{(a) Linear approximation: $dx$ and $dy$ as numbers}

\begin{signature}
  $dx \in \R$; \quad $dy = f'(x)\,dx \in \R$.
\end{signature}

Here $dx$ is not a function of $x$ at all: it is an independent real number, freely chosen, playing the role of an input increment. Given $x$ and $dx$, the differential $dy$ is the value of the map
\[
  \mathrm{d}f : I \times \R \to \R, \qquad (x, dx) \mapsto f'(x)\,dx,
\]
which is linear in its second argument for each fixed $x$. Ordinary notation shows only ``$dx$'' on the page and lets the dependence on $x$ stay implicit in ``the point where we're differentiating''---but $dy$ genuinely depends on \emph{two} numbers, not one.

\begin{center}
  \renewcommand{\arraystretch}{1.35}
  \begin{tabular}{@{}l l l l@{}}
    \toprule
    \textbf{Symbol} & \textbf{Type} & \textbf{Domain / Codomain} & \textbf{Key property} \\
    \midrule
    $x$  & point   & $x \in I$                     & fixed base point \\
    $dx$ & number  & $dx \in \R$, arbitrary, $\neq 0$ & independent variable \\
    $dy$ & number  & $dy \in \R$                    & $dy = f'(x)\,dx$; linear in $dx$ \\
    \bottomrule
  \end{tabular}
\end{center}

\subsection*{(b) Linear map: $dx$ and $dy$ as functionals on tangent vectors}

\begin{signature}
  $dx_p, dy_p : T_p\R^2 \to \R$; \quad $dF_p = \dfrac{\partial F}{\partial x}(p)\,dx_p + \dfrac{\partial F}{\partial y}(p)\,dy_p : T_p\R^2 \to \R$.
\end{signature}

Fix a point $p = (x,y) \in \R^2$ and its tangent space $T_p\R^2 \cong \R^2$, whose elements are vectors $(h,k)$. Here $dx_p$ and $dy_p$ are not numbers but \emph{linear functionals}---elements of the cotangent space $T_p^*\R^2$---namely the coordinate projections $dx_p(h,k) = h$ and $dy_p(h,k) = k$. The partial derivatives $\partial F/\partial x, \partial F/\partial y : \R^2 \to \R$ are ordinary scalar-valued functions of the point $p$ (they eat a point, return a number), whereas $dF_p$ is again a linear functional on $T_p\R^2$, built as the linear combination above. Letting $p$ vary, $dF$ is a \emph{section of the cotangent bundle} $T^*\R^2$---a $1$-form, i.e.\ an assignment of a functional to every point, not a single functional or a single number. The same goes for $M, N : \R^2 \to \R$ (scalar functions of $p$) versus $\omega = M\,dx + N\,dy$ (a $1$-form, $\omega_p = M(p)\,dx_p + N(p)\,dy_p$).

\begin{center}
  \small
  \renewcommand{\arraystretch}{1.35}
  \begin{tabular}{@{}l l l l@{}}
    \toprule
    \textbf{Symbol} & \textbf{Type} & \textbf{Domain / Codomain} & \textbf{Key property} \\
    \midrule
    $p=(x,y)$ & point            & $p \in \R^2$                        & evaluation point \\
    $\partial F/\partial x$ & scalar field & $\R^2 \to \R$                 & function of $p$ alone \\
    $dx_p, dy_p$ & linear functional & $T_p\R^2 \to \R$                & coordinate projections \\
    $dF_p$    & linear functional & $T_p\R^2 \to \R$                   & $\tfrac{\partial F}{\partial x}(p)\,dx_p + \tfrac{\partial F}{\partial y}(p)\,dy_p$ \\
    $dF$, $\omega$ & $1$-form (section) & $p \mapsto (T_p\R^2 \to \R)$ & functional at \emph{every} point \\
    \bottomrule
  \end{tabular}
\end{center}

\subsection*{(c) Infinitesimal: $dx$ and $dy$ as hyperreal numbers}

\begin{signature}
  $dx \in \R^*$ infinitesimal; \quad $dy = {}^*\!f(x+dx) - {}^*\!f(x) \in \R^*$; \quad $\operatorname{st}(dy/dx) = f'(x) \in \R$.
\end{signature}

Let $\R^*$ be the hyperreal field extending $\R$. Here $dx \in \R^*$ is a genuine, fixed, nonzero number---just an infinitesimal one, meaning $\lvert dx\rvert < r$ for every positive real $r$. The function $f$ extends, via the transfer principle, to ${}^*\!f : \R^* \to \R^*$, and $dy$ is \emph{defined} as the true increment $dy := {}^*\!f(x+dx) - {}^*\!f(x) \in \R^*$, itself infinitesimal when $f$ is continuous at $x$. The quotient $dy/dx$ is an ordinary quotient of hyperreal numbers, living in $\R^*$; generically it is \emph{finite} (not infinitesimal), which is what allows the standard-part map $\operatorname{st}$, defined on finite hyperreals, $\operatorname{st} : \mathrm{fin}(\R^*) \to \R$, to send it back to an ordinary real number equal to $f'(x)$.

\begin{center}
  \small
  \renewcommand{\arraystretch}{1.35}
  \begin{tabular}{@{}l l l l@{}}
    \toprule
    \textbf{Symbol} & \textbf{Type} & \textbf{Domain / Codomain} & \textbf{Key property} \\
    \midrule
    $dx$        & hyperreal number & $dx \in \R^*$, infinitesimal $\neq 0$ & $\lvert dx\rvert < r$, all real $r>0$ \\
    ${}^*\!f$   & function        & $\R^* \to \R^*$                       & transfer of $f$ \\
    $dy$        & hyperreal number & $dy \in \R^*$                        & $:= {}^*\!f(x+dx) - {}^*\!f(x)$ \\
    $dy/dx$     & hyperreal number & $\in \mathrm{fin}(\R^*) \subset \R^*$ & finite, not infinitesimal \\
    $\operatorname{st}$ & function & $\mathrm{fin}(\R^*) \to \R$          & $\operatorname{st}(dy/dx) = f'(x)$ \\
    \bottomrule
  \end{tabular}
\end{center}

\subsection*{Same notation, different type}

The punchline of this section is visible only by comparing rows: three frameworks reuse the letters $dx$ and $dy$ for objects that do not even live in the same kind of set.

\begin{center}
  \renewcommand{\arraystretch}{1.35}
  \begin{tabular}{@{}l l l@{}}
    \toprule
    \textbf{Framework} & \textbf{Type of $dx$} & \textbf{Type of $dy$} \\
    \midrule
    (a) Linear approximation & number, $dx \in \R$                & number, $dy \in \R$ \\
    (b) Linear map            & functional, $dx_p \in T_p^*\R^2$   & functional, $dy_p \in T_p^*\R^2$ \\
    (c) Infinitesimal         & hyperreal number, $dx \in \R^*$    & hyperreal number, $dy \in \R^*$ \\
    \bottomrule
  \end{tabular}
\end{center}

Only within framework (c) is $dx$ literally a nonzero number that $dy$ is divided by to form a bona fide quotient; in (a) the ``fraction'' $dy/dx$ is trivial cancellation by construction; in (b) there is no division at all, since $dx_p$ and $dy_p$ are functionals, not numbers, and $dF/dx$ is recovered by \emph{evaluating} $dF_p$ on a tangent vector, not by dividing.

\section{Why Exact Equations Are Best Written as $M\,dx + N\,dy = 0$}

Framing (b) reveals the deeper reason for this notation, beyond mere convenience. Regard
\[
\omega = M\,dx + N\,dy
\]
as a \emph{differential $1$-form}: at each point $(x,y)$ it is the linear functional $\omega_{(x,y)}(h,k) = M(x,y)\,h + N(x,y)\,k$ acting on tangent vectors. The equation $\omega = 0$ asserts that along a solution curve with velocity $(x'(t), y'(t))$,
\[
M\,x'(t) + N\,y'(t) = 0.
\]
Parametrizing the curve by $x$ itself, so that $t = x$, this becomes $M + N\,dy/dx = 0$---the original ODE. The ``cross-multiplication'' is thus not a trick applied to a fraction; it is the coordinate-free statement $\omega = 0$, read off in the $x$-parametrization. Writing $M\,dx + N\,dy = 0$ instead of $dy/dx = -M/N$ is deliberately symmetric in $x$ and $y$: it does not presuppose which variable is a function of the other, which matters when one later solves for $x(y)$ instead of $y(x)$.

This symmetry is also why the exactness criterion,
\[
\frac{\partial M}{\partial y} = \frac{\partial N}{\partial x},
\]
works as it does: it is precisely the statement that $\omega$ is a \emph{closed} form. By the Poincar\'e lemma, a closed $1$-form on a simply connected domain is \emph{exact}, meaning $\omega = dF$ for some potential $F$, in exactly the sense of (b) above. The name ``exact differential equation'' is therefore not a metaphor---it invokes the exact/closed distinction for differential forms directly.

\section{Conclusion}

Elementary calculus is right that the derivative, defined as a limit, is not a fraction. But the differential is a distinct, later-introduced notion: a real number (or a linear functional, or an infinitesimal, depending on the chosen framework) satisfying $dy = f'(x)\,dx$ by construction rather than by limiting process. Once this relationship is taken as a definition, dividing or cross-multiplying it is ordinary algebra. The apparent contradiction dissolves once one recognizes that calculus reuses a single piece of notation, $dy/dx$, for two related but logically distinct objects, and that the switch between them is rarely announced.

What Section~5 adds is that this is not merely a two-way ambiguity between ``derivative'' and ``differential.'' Even within ``the differential,'' the type of $dx$ and $dy$---number, linear functional, or hyperreal number---depends on which of the three frameworks is in force. Writing down the equation $dy = f'(x)\,dx$ is not enough to say what $dx$ \emph{is}; only its domain and codomain, stated explicitly, do that.

\appendix

\section{Linear Maps and Tangent Spaces}
\label{app:background}

Framework (b) in Section~5 leans on two ideas that elementary calculus does not usually name explicitly: \emph{linear maps} and \emph{tangent spaces}. Both are simpler than their names suggest.

\subsection*{Linear maps}

\begin{definition}
A function $L : \R^2 \to \R$ is a \emph{linear map} (or \emph{linear functional}) if, for all $(h_1,k_1), (h_2,k_2) \in \R^2$ and all scalars $c \in \R$,
\[
L(h_1+h_2,\,k_1+k_2) = L(h_1,k_1) + L(h_2,k_2), \qquad L(ch_1,ck_1) = c\,L(h_1,k_1).
\]
\end{definition}

In down-to-earth terms, a linear map out of $\R^2$ is nothing more exotic than an expression
\[
L(h,k) = a\,h + b\,k
\]
for fixed constants $a, b \in \R$---and every such expression is automatically linear, so this formula is not a special case but the \emph{general} one. Contrast it with an affine function like $L(h,k) = ah+bk+c$ with $c \ne 0$: adding a constant term destroys linearity, since a linear map must send $(0,0) \mapsto 0$, while $L(0,0) = c \ne 0$ here.

\begin{example}
$L(h,k) = 3h - 2k$ is linear: $L(1,0)=3$, $L(0,1)=-2$, and $L(2,5) = 6-10=-4$, which agrees with $2L(1,0)+5L(0,1) = 6-10=-4$, as linearity requires.
\end{example}

The two simplest linear maps on $\R^2$ just read off a coordinate: $\pi_1(h,k)=h$ and $\pi_2(h,k)=k$. These are exactly $dx_p$ and $dy_p$ from Section~5(b): as \emph{formulas} they do not even mention $p$, which is why the subscript is there only to record \emph{which} tangent space they are regarded as acting on (see below). Every other linear map on $\R^2$, such as $dF_p(h,k) = F_x(p)\,h+F_y(p)\,k$, is a constant-coefficient combination $a\pi_1+b\pi_2$ of these two---this is what Section~5(b) means by expressing $dF$ ``in the dual basis $\{dx,dy\}$.''

\subsection*{Tangent vectors and the tangent space}

Fix a point $p=(x_0,y_0)\in\R^2$ and a curve through $p$: differentiable functions $x(t),y(t)$ with $x(t_0)=x_0$, $y(t_0)=y_0$ at some parameter value $t_0$. Its \emph{velocity vector} at $p$ is
\[
(x'(t_0),\,y'(t_0)) \in \R^2,
\]
the instantaneous direction and speed of motion through $p$. Different curves through $p$ can share the same velocity vector, so the map from curves to velocity vectors is many-to-one. Conversely, every pair $(h,k)\in\R^2$ arises as the velocity vector of some curve through $p$---for instance, the straight line $x(t)=x_0+th$, $y(t)=y_0+tk$ has velocity $(h,k)$ at $t_0=0$. Hence, ranging over all curves through $p$, the set of attainable velocity vectors is all of $\R^2$.

\begin{definition}
The \emph{tangent space} to $\R^2$ at $p$, written $T_p\R^2$, is the set of velocity vectors of curves through $p$. As a set it is identical to $\R^2$, but its elements are pictured as \emph{arrows based at $p$}---directions and speeds of motion---rather than as points or positions.
\end{definition}

Why distinguish $T_p\R^2$ from $\R^2$ if they are the same set? Because ``a location'' and ``a direction of travel through that location'' are conceptually different, even though both happen to be pairs of numbers in the flat plane. The distinction is invisible in $\R^2$ precisely because the plane is flat: on a curved surface (a sphere, say) the tangent space at a point is a flat plane genuinely different from the surface itself, and tangent spaces at two different points cannot be identified with each other. Framework (b) writes $T_p\R^2$, not just $\R^2$, so that its notation carries over unchanged to that curved setting---though for this essay $T_p\R^2$ can always be read simply as ``$\R^2$, but thought of as velocity vectors at $p$.''

A linear functional on the tangent space---an element of $T_p^*\R^2$---is a linear map that eats a tangent vector $(h,k)\in T_p\R^2$ and returns a number. This is precisely the type of $dx_p$, $dy_p$, and $dF_p$ in Section~5(b): each turns a direction of travel through $p$ into a single real number.

\begin{example}
Let $p=(1,1)$ and take the curve $x(t)=t$, $y(t)=t^2$, through $p$ at $t=1$ with velocity $(x'(1),y'(1)) = (1,2) \in T_p\R^2$. For $F(x,y)=x^2y$, $F_x(p) = 2xy|_p = 2$ and $F_y(p) = x^2|_p=1$, so
\[
dF_p(1,2) = F_x(p)\cdot 1 + F_y(p)\cdot 2 = 2\cdot1+1\cdot2 = 4.
\]
This matches the ordinary chain rule directly: $\frac{d}{dt}F(x(t),y(t))\big|_{t=1} = \frac{d}{dt}(t^4)\big|_{t=1} = 4t^3\big|_{t=1}=4$. Evaluating the linear functional $dF_p$ on the tangent vector $(1,2)$ \emph{is} evaluating this derivative---the general fact Section~5(b) uses to recover the chain rule without splitting a fraction.
\end{example}

\section*{Self-Check Questions}

\begin{enumerate}[leftmargin=*, itemsep=4pt]
  \item In the identity $\dfrac{dF}{dx} = \dfrac{\partial F}{\partial x} + \dfrac{\partial F}{\partial y}\dfrac{dy}{dx}$, is $dy/dx$ being split into two fractions? Justify your answer.
  \item State the precise definition of $dy/dx$ as a derivative (the limit definition), and explain why this object cannot literally be a quotient of two numbers.
  \item State the definition of the differential $dy$ in terms of $dx$, once $f'(x)$ is known. Why does dividing this equation by $dx$ trivially recover $f'(x)$?
  \item Explain, in your own words, why the derivative and the differential can share the same notation $dy/dx$ without being the same kind of object.
  \item Under the linear-approximation view of the differential, is $dy$ equal to $\Delta y$? What is the precise relationship between them?
  \item In the linear-map definition, what are $dx$ and $dy$, formally? What role do they play as elements of a dual basis?
  \item Using the linear-map definition, derive the chain-rule identity $dF/dx = \partial F/\partial x + (\partial F/\partial y)(dy/dx)$ without invoking fractions at any point.
  \item What is an infinitesimal hyperreal, and how does the non-standard analysis definition of $dy/dx$ differ from the limit definition? What operation recovers $f'(x)$ from the hyperreal quotient?
  \item Why is $M\,dx + N\,dy = 0$ considered a more natural way to state an ODE than $dy/dx = -M/N$? What symmetry does it preserve?
  \item Explain the geometric meaning of a differential $1$-form $\omega = M\,dx + N\,dy$, and what it means for $\omega$ to vanish along a curve.
  \item State the exactness condition $\partial M/\partial y = \partial N/\partial x$ and explain, using the Poincar\'e lemma, why it guarantees the existence of a potential function $F$ with $dF = \omega$.
  \item Why is the term ``exact'' in ``exact differential equation'' not merely descriptive but a direct reference to the exact/closed distinction for differential forms?
  \item Suppose someone tells you ``$dy/dx$ is just a fraction, you can always cross-multiply.'' Identify precisely where this claim is true and where it requires qualification.
  \item In framework (a), is $dy$ a function of one variable or two? Name them, and explain why ordinary notation only displays one.
  \item In framework (b), is $dx_p$ a number or a function? State its domain and its codomain precisely.
  \item Still in framework (b), what is the difference in type between $\partial F/\partial x$ and $dF_p$? What is the difference in type between $dF_p$ (fixed $p$) and $dF$ (as $p$ varies)?
  \item In framework (c), what is the codomain of the standard-part map $\operatorname{st}$, and why must $dy/dx$ be \emph{finite} (as opposed to infinite or infinitesimal) for $\operatorname{st}(dy/dx)$ to be defined?
  \item Compare the three frameworks' treatment of $dx$: in which one, if any, is $dx$ literally a nonzero number that $dy$ can be divided by to form a genuine quotient? What replaces division in the other framework(s)?
  \item Explain why stating the equation $dy = f'(x)\,dx$ alone does not pin down what kind of object $dx$ is, and why the domain/codomain of $dx$ must be specified separately.
  \item Give the general formula for a linear map $L : \R^2 \to \R$, and explain why an affine function $L(h,k) = ah+bk+c$ with $c \ne 0$ fails to be linear.
  \item What is a tangent vector at a point $p \in \R^2$, and why is the tangent space $T_p\R^2$ conceptually different from $\R^2$ even though the two sets coincide?
\end{enumerate}

\end{document}
